\documentclass{umk-matematika}
\begin{document}

\begin{center}
{\Large\bfseries Практическая работа 20}
\end{center}
\vspace{0.5cm}

\textbf{Задача 1.} Измерения прямоугольного параллелепипеда: $3$ см, $4$ см, $5$ см. Найдите его объём.\par\vspace{10pt}
\textbf{Задача 2.} Объём прямоугольного параллелепипеда равен $120$ см³. Два измерения равны $4$ см и $5$ см. Найдите третье измерение.\par\vspace{10pt}
\textbf{Задача 3.} Куб имеет ребро $4$ см. Найдите его объём.\par\vspace{10pt}
\textbf{Задача 4.} Площадь основания прямоугольного параллелепипеда $15$ см², высота $7$ см. Найдите объём.\par\vspace{10pt}
\textbf{Задача 5.} Диагональ прямоугольного параллелепипеда равна $13$ см. Измерения основания $3$ см и $4$ см. Найдите объём.\par\vspace{10pt}
\textbf{Задача 6.} Объём куба равен $125$ см³. Найдите площадь его полной поверхности.\par\vspace{10pt}
\textbf{Задача 7.} В прямоугольном параллелепипеде стороны основания $6$ см и $8$ см. Диагональ параллелепипеда наклонена к плоскости основания под углом $45^\circ$. Найдите объём.\par\vspace{10pt}
\textbf{Задача 8.} Фундамент здания — прямоугольный параллелепипед с размерами $12$ м $\times$ $8$ м $\times$ $1{,}5$ м. Внутри вырезан прямоугольный проём для подвала размером $6$ м $\times$ $4$ м на всю высоту фундамента. Найдите объём бетона.\par\vspace{10pt}
\textbf{Задача 9.} Бак для воды имеет форму прямоугольного параллелепипеда с квадратным основанием. Объём бака $1000$ литров, высота $1$ м. Найдите сторону основания в метрах. (1 м³ = 1000 л)\par\vspace{10pt}
\textbf{Задача 10.} Комната имеет форму прямоугольного параллелепипеда. Площадь пола $20$ м², площадь одной стены $15$ м², площадь смежной стены $12$ м². Найдите объём комнаты.\par\vspace{10pt}
\newpage
\section*{Ответы}

\begin{enumerate}
  \item Ответ: $60$ см³
  \item Ответ: $6$ см
  \item Ответ: $64$ см³
  \item Ответ: $105$ см³
  \item Ответ: $144$ см³
  \item Ответ: $150$ см²
  \item Ответ: $480$ см³
  \item Ответ: $108$ м³
  \item Ответ: $1$ м
  \item Ответ: $60$ м³
\end{enumerate}
\end{document}
